\chapter{CLI Surface and End-to-End Data Flow}\label{ch:cli-dataflow}

% Local helpers for TikZ node text, where \code (url's \path) cannot be used
% because TikZ redefines \path inside a tikzpicture.
\providecommand{\dftt}[1]{\texttt{#1}}

\lead{What this chapter covers}
This chapter follows a command line from the console script to the last file it writes. The stage chapters that follow (Chapters~\ref{ch:s01-s02} to \ref{ch:s13-s14}) hang on this map. It fixes the stage identifiers \leavevmode\stagetag{S01} to \leavevmode\stagetag{S14} in execution order. It names the function that runs at each step, with its line in the working tree of 5 November 2026, and lists the files each step writes. It also records where the README and other documents describe a different program from the one the code runs. The markdown twin of this chapter is \code{docs/architecture/DATAFLOW.md}.

\lead{Basis and notation}
Everything below comes from a code-derived trace of the working tree. The trace combines an abstract-syntax-tree (AST) import closure with a read of the dispatch code, and it was spot-checked line by line against \code{raosim/run_nozzle.py}. That file has 6,981 lines, and its \code{main()} alone spans lines 2817 to 6977. A bare line reference such as \code{:4452} means a line of \code{raosim/run_nozzle.py}; references into other files name the file. Where a document and the code disagree, the code is authoritative and the disagreement is listed in Section~\ref{sec:df-drift}.

\caution{Every path in this chapter produces a preliminary, screening-level design product. The runner sets \code{summary["hardware_qualified"]=False} unconditionally (\code{:4968}) and ends every successful run by printing that numerical design and CAD output is never hardware-qualified by LREKit (\code{:6976}). The job run record writes \code{"hardware_release":False} into its manifest (\code{raosim/mdo/run_record.py:46}). Nothing in this data flow is physically validated.}

% ==============================================================================
\section{Entry points}\label{sec:df-entry}

\code{pyproject.toml} declares four console scripts (lines 62 to 66) and packages the top-level \code{main.py} as a module (line 69). Together with \code{python}~\code{-m} and two compatibility wrappers, these reach three dispatch functions.

{\small
\begin{longtable}{P{3.5cm}P{4.2cm}P{7.8cm}}
\toprule
Entry & Target & What runs \\
\midrule
\endhead
\code{lrekit} & \code{lrekit.cli:main} (\code{lrekit/cli.py:18}) & If \code{argv[1]} is one of the \code{JOB_VERBS} \code{guided}, \code{run}, \code{validate} or \code{catalog} (\code{lrekit/job_cli.py:29}), control passes to \code{lrekit.job_cli.main} (\code{job_cli.py:398}). Otherwise control passes to \code{raosim.cli.main}. \\
\code{raosim}, \code{RaoRocketSim} & \code{raosim.cli:main} (\code{raosim/cli.py:9}) & If \code{argv[1]} is \code{export-openfoam-spray} (\code{raosim/cli.py:11}), control passes to \code{raosim.openfoam.case.cli_main} (\code{openfoam/case.py:1611}). Otherwise the runner \code{raosim.run_nozzle.main} (\code{:2817}) is imported lazily and run. \\
\code{RaoRocketSimLegacy} & \code{raosim.cli:legacy_main} (\code{raosim/cli.py:21}) & \code{import_module("main").main()}, which is the old top-level \code{main.py} toolbox (cross-cutting group X03). \\
\code{python}~\code{-m}~\code{lrekit}, \code{python}~\code{-m}~\code{raosim} & \code{lrekit/__main__.py}, \code{raosim/__main__.py} & Same as \code{lrekit} and \code{raosim}. \\
\code{scripts/run_nozzle.py}, \code{raosim/run.sh} & compatibility wrappers & \code{raosim.run_nozzle.main}. \\
\bottomrule
\end{longtable}}

The two package names are tied together by one line. \code{lrekit/__init__.py:16} sets \code{__path__} to the \code{lrekit/} directory followed by \code{raosim.__path__}, so \code{lrekit.<x>} resolves to \code{raosim/<x>.py}. There is no second copy of any engineering module.

Dispatch is lazy. Eagerly, \code{lrekit.cli} imports only four modules: \code{lrekit/__init__.py}, \code{lrekit/cli.py}, \code{raosim/__init__.py} and \code{raosim/cli.py}. Everything else loads inside functions. Importing \code{raosim/run_nozzle.py} eagerly pulls in 19 modules: \code{chamber_geometry}, \code{coolants}, \code{export}, \code{gas_dynamics}, \code{materials}, \code{moc}, \code{nasa_moc}, \code{physics}, \code{rao_residuals}, \code{rao_variational}, \code{regen_geometry}, \code{regen_profile}, \code{separation}, \code{thermofluids}, \code{throat_geometry}, \code{validation}, \code{wall_model}, the \code{raosim} package itself and the runner. The disciplines that only some flags need, such as injector, pumps, CAD and MDO, load on first use.

% ==============================================================================
\section{CLI surface}\label{sec:df-surface}

\subsection{Job verbs and the OpenFOAM export verb}\label{sec:df-verbs}

Five verbs bypass the flag runner. Four of them belong to the job dialect: YAML or JSON job documents with unit-bearing values, compiled into the versioned \code{EngineJobSpec} of \code{raosim/mdo/job.py}. The fifth writes an OpenFOAM case.

{\small
\begin{longtable}{P{4.6cm}P{6.0cm}P{4.9cm}}
\toprule
Invocation & Handler chain & Writes \\
\midrule
\endhead
\code{lrekit}~\code{guided} \code{[--save}~\code{FILE]} \code{[--run]} & \code{job_cli._cmd_guided} (\code{:274}) $\to$ \code{lrekit.guided.run_guided} (\code{guided.py:271}). That calls \code{GuidedSession.run}, an interview that re-prompts through \code{units.parse_quantity}. Then \code{job_cli._build_spec}, \code{mdo.job.resolve_engine_job} (\code{job.py:889}) and a coverage preview through \code{mdo.coolant_htd.htd_availability}. With \code{--run}, \code{_cmd_run} follows. & Job YAML or JSON with \code{--save}. \\
\code{lrekit}~\code{validate}~\code{JOB.yaml} \code{[--resolved-out}~\code{FILE]} & \code{_cmd_validate} (\code{job_cli.py:181}): \code{_load_job_document} (rejects duplicate keys), \code{_build_spec} (rejects unknown keys), \code{mdo.execute.expand_architectures} (\code{execute.py:814}) and \code{resolve_engine_job}. Nothing is solved. & Resolved job JSON with \code{--resolved-out}. \\
\code{lrekit}~\code{run}~\code{JOB.yaml} \code{[--out}~\code{FILE]} & \code{_cmd_run} (\code{job_cli.py:209}). Section~\ref{sec:df-job} gives the full path. & Run record under \code{output_policy.output_dir}; result JSON with \code{--out}. The exit code is the verdict. \\
\code{lrekit}~\code{catalog}~\code{TOPIC} & \code{_cmd_catalog} (\code{job_cli.py:320}). The eight topics (\code{job_cli.py:31-34}) are \code{variables}, \code{targets}, \code{objectives}, \code{constraints}, \code{propellants}, \code{materials}, \code{unavailable} and \code{capabilities}. They read \code{mdo.job.VARIABLE_DIMENSIONS}, \code{mdo.metrics}, \code{mdo.constraints.CONSTRAINT_MANIFEST}, \code{mdo.propellants.PROPELLANTS}, \code{materials.material_names} and \code{mdo.capabilities.JOB_CAPABILITIES}. & Console only. \\
\code{lrekit}~\code{export-openfoam-spray} \code{OUTPUT} \code{[--case-row}~\code{case_1..case_4]} \code{[--wedge-angle-deg]} \code{[--opening-cells]} \code{[--radial-cells]} & \code{raosim/cli.py:11-13} $\to$ \code{openfoam.case.cli_main} (\code{case.py:1611}) $\to$ \code{build_radhakrishnan2018_sheet_vof_case} (\code{:456}) $\to$ \code{write_openfoam_case} (\code{:1520}). The defaults are \code{case_1}, $5.0^\circ$, 6 and 240 (\code{case.py:1622-1627}). & A version-pinned OpenFOAM volume-of-fluid (VOF) case plus \code{raosim_openfoam_manifest.json} (\code{case.py:47}). OpenFOAM itself is never run. \\
\bottomrule
\end{longtable}}

\subsection{The flag runner: 330 flags by family}\label{sec:df-flags}

\code{raosim.run_nozzle.main} builds one flat \code{argparse.ArgumentParser} with 330 \code{add_argument} calls (lines 2819 to 3782). It uses no argument groups. The grouped list that \code{--tags} prints comes from \code{print_tags} (\code{:139}). That list is maintained by hand, separately from the parser, and shows only a subset of the flags. Table~\ref{tab:df-flag-families} sorts all 330 flags into families by the stage that consumes them. The grouping is this guide's own. The counts come from an AST walk of \code{main()} and sum to 330.

{\small
\begin{longtable}{P{3.3cm}P{0.6cm}P{8.2cm}P{2.6cm}}
\caption{Runner flag families (AST count of \texttt{add\_argument} calls in \texttt{run\_nozzle.main}, working tree 2026-11-05).}\label{tab:df-flag-families}\\
\toprule
Family & $n$ & Representative flags & Consumed by \\
\midrule
\endfirsthead
\toprule
Family & $n$ & Representative flags & Consumed by \\
\midrule
\endhead
Operating point and nozzle & 9 & \code{--pc} \code{--pa-over-p0} \code{--target-thrust} \code{--rt} \code{--epsilon} \code{--length-pct} \code{--gamma} \code{--contour-method} \code{--allow-chart-extrapolation} & \leavevmode\stagetag{S04}, \leavevmode\stagetag{S05} \\
Propellant and thermochemistry & 15 & \code{--propellant} \code{--oxidizer} \code{--fuel} \code{--thermo-mode} \code{--mixture-ratio} \code{--eta-cstar} \code{--eta-cf} \code{--nozzle-expansion-model} \code{--frozen-gas-table}, plus the per-stream inlet temperature, pressure and phase flags & \leavevmode\stagetag{S03}; inlet states \leavevmode\stagetag{S08} \\
Chamber and throat geometry & 11 & \code{--l-star} \code{--contraction-ratio} \code{--shoulder-radius-factor} \code{--shoulder-sizing} \code{--minimum-cylindrical-length} \code{--convergent-angle} \code{--ru-factor} \code{--rd-factor} \code{--cd-target} & \leavevmode\stagetag{S05} \\
Rao BVP solver & 8 & \code{--backend} \code{--max-nfev} \code{--n-control} \code{--n-kernel} \code{--theta-b-guess} \code{--kernel-dtheta-limit-deg} \code{--kernel-massflow-rel-tol} \code{--allow-unconverged} & \leavevmode\stagetag{S06} only \\
Release evidence & 3 & \code{--release-evidence-manifest} \code{--configuration-id} \code{--require-release-evidence} & \leavevmode\stagetag{S14} \\
Regen channels, cooling, thermofluids, wall sizing and structure & 47 & \code{--regen} \code{--channels} \code{--channel-*} \code{--helix-turns} \code{--wall-thickness} \code{--wall-sizing} \code{--size-wall} \code{--auto-size} \code{--thermal} \code{--material} \code{--jacket-material} \code{--coolant*} \code{--coolant-property-backend} \code{--hydraulic-network} \code{--radiation-*} \code{--boiling-chf} \code{--structural-fos} \code{--required-cycles} \code{--t-hot-*} \code{--buckling-*} & \leavevmode\stagetag{S07}; regen geometry \leavevmode\stagetag{S11} \\
Fixed pintle injector, feed ports, manifolds and throttle map & 47 & \code{--injector} \code{--injector-architecture} \code{--injector-sizing} \code{--fuel-injector-dp-fraction} \code{--oxidizer-injector-dp-fraction} \code{--pintle-*} \code{--*-discharge-coefficient} \code{--*-inlet-count/diameter/angle/fitting} \code{--*-manifold-width/depth} \code{--allow-infeasible-injector} \code{--throttle-map}. Also the hidden \code{--injector-pressure-drop}, whose help is suppressed and which is rejected at parse time (\code{:661}) & \leavevmode\stagetag{S08}, \leavevmode\stagetag{S10} \\
Movable (Son) pintle evidence & 54 & \code{--movable-post-*} \code{--movable-center-gap-*} \code{--movable-cd-*} (calibration map, provenance and validity ranges) \code{--movable-position-*} \code{--movable-actuator-*} \code{--movable-stem-*} \code{--movable-sheet-*} & \leavevmode\stagetag{S08} \\
Spray/$c^*$ coupling & 7 & \code{--spray-cstar-coupling} \code{--spray-eta-mixing} \code{--spray-eta-combustion} \code{--spray-coupling-*} \code{--spray-evaporation-constant} & \leavevmode\stagetag{S08} \\
Feed-pressure ledger & 23 & \code{--feed-architecture} \code{--*-tank-pressure} \code{--*-supply-pressure} \code{--*-flow-capacity} \code{--*-line-loss[-fraction]} \code{--*-manifold-loss[-fraction]} \code{--*-control-margin[-fraction]} \code{--*-npsh-required} \code{--pump-efficiency-*} & \leavevmode\stagetag{S08} feed spec, \leavevmode\stagetag{S09} \\
Electric pump, motor, inverter and battery & 27 & \code{--electric-pump} (with \code{--no-electric-pump}) \code{--electric-technology} \code{--pump-rpm} \code{--burn-time} \code{--motor-*} \code{--inverter-*} \code{--battery-*} \code{--vehicle-mass} \code{--pump-head-coefficient} \code{--pump-tip-speed-limit} \code{--allow-infeasible-pump} & \leavevmode\stagetag{S09} \\
Interface, seals and manufacturing limits & 19 & \code{--flange-od} \code{--flange-length} \code{--bolt-*} \code{--joint-separation-factor} \code{--injector-face-od} \code{--injector-face-thickness} \code{--seal-type} \code{--o-ring-groove-*} \code{--min-tool-diameter} \code{--injector-tolerance} \code{--igniter-port-*} & \leavevmode\stagetag{S08}, \leavevmode\stagetag{S10}, \leavevmode\stagetag{S11} \\
CAD export & 16 & \code{--cad} \code{--require-brep} \code{--regen-brep} \code{--regen-manifolds} \code{--regen-release-mode} \code{--regen-port-*} \code{--engine-assembly} \code{--injector-cad} \code{--injector-cad-format} \code{--pintle-sleeve} \code{--pump-cad} \code{--pump-cad-format} \code{--allow-open-pump-mesh} & \leavevmode\stagetag{S11} \\
Visualisation & 5 & \code{--flowfield} \code{--animate} \code{--chamber-temp} \code{--pump-visualize} \code{--show} & \leavevmode\stagetag{S12} \\
Modes and output & 6 & \code{-i/--interactive} \code{--complete-package} \code{--tags} \code{--list-materials} \code{--no-banner} \code{--out} & \leavevmode\stagetag{S01} \\
Whole-engine MDO & 23 & \code{--engine-mdo} \code{--engine-mdo-optimize} \code{--engine-mdo-couple-cstar} \code{--isp-min} \code{--isp-sweep} \code{--mdo-multistart} \code{--mdo-station-study} \code{--mdo-propellant} \code{--mdo-of-mode} \code{--mdo-export} \code{--design-margins} \code{--film-frac} \code{--t-wall} \code{--allow-incomplete-physics} & \leavevmode\stagetag{S13}, \leavevmode\stagetag{S14} \\
Requirement-driven design & 10 & \code{--requirements} \code{--thrust-condition} \code{--isp-basis} \code{--flight-duration} \code{--qualification-duration} \code{--envelope-diameter-max} \code{--envelope-length-max} \code{--burnout-mass-max} \code{--throttle-range} \code{--reusable-cycles} & \leavevmode\stagetag{S13} \\
\midrule
Total & 330 & & \\
\bottomrule
\end{longtable}}

Three behaviours of the parser matter for every run:
\begin{itemize}
  \item \textbf{Argument files.} \code{_expand_arg_files} (\code{:562}, called at \code{:3785}) expands \code{@file.args} tokens before parsing. Each file is read as shell-style tokens with comments, nested includes resolve relative to the including file, and expansion is positional, so a later flag overrides an earlier one. The four shipped files live in \code{examples/cli/}.
  \item \textbf{Units.} Every runner flag is an \code{argparse} number read as SI by convention (Pa, m, N, kg/s). The runner never calls \code{raosim/units.py}; only the job dialect parses unit-bearing strings.
  \item \textbf{Defaults that decide the path.} \code{--contour-method} defaults to \code{bezier} (\code{:2827-2831}), \code{--out} to \code{builds/nozzle_run} (\code{:3534}), \code{--cad} to \code{none} (\code{:3535}) and \code{--require-brep} to true (\code{:3538-3545}). \code{--backend} defaults to \code{jax} (\code{:2953}), but that default only matters on the \code{rao-bvp} path.
\end{itemize}

\caution{The eight Rao BVP solver flags act only when \code{--contour-method}~\code{rao-bvp} is given. On the default B\'ezier path, \code{_solve} (\code{:1464-1487}) calls \code{nozzle_geometry.bell_nozzle_contour} with none of them, so \code{--backend}, \code{--max-nfev}, \code{--n-control} and \code{--n-kernel} have no effect. Three of the four shipped argument files pass these flags without selecting \code{rao-bvp} (Section~\ref{sec:df-drift}, item 9).}

% ==============================================================================
\section{Dispatch inside the runner}\label{sec:df-dispatch}

Before any physics runs, \code{run_nozzle.main} applies these steps in this order:
\begin{enumerate}
  \item It builds the parser (\code{:2819-3782}), expands argument files (\code{:3785}) and parses (\code{:3790}).
  \item \code{_reject_legacy_injector_pressure_drop} (\code{:661}, called at \code{:3791}) exits through \code{argparse} if the deprecated \code{--injector-pressure-drop} appears.
  \item \code{_apply_electric_technology} (\code{:522}, called at \code{:3792}) resolves \code{raosim.electric_technology.get_electric_technology} and fills the eight motor, inverter and battery flags.
  \item \textbf{Front door} (\code{:3797-3811}). This step runs on a bare invocation or with \code{-i}, unless \code{--engine-mdo} or \code{--engine-mdo-optimize} is set. \code{_select_mode_interactively} (\code{:1850}) prints \code{_MODE_MENU} (\code{:1756}). Choice~1 continues on the traditional path with the starter defaults. Choice~2 runs \code{_interactive_engine_mdo} (\code{:1778}) and then \code{_run_engine_mdo}. Choice~3 runs the same interview with \code{optimise=True} and then \code{_run_engine_mdo_optimize}. End of input returns 130.
  \item \code{_validate_common_engine_args} (\code{:719}, called at \code{:3812}) calls \code{raosim.input_validation.validate_engine_inputs}. This is the shared domain check, and it runs before every branch.
  \item \code{--mdo-multistart}, \code{--mdo-multistart-seed} and \code{--mdo-station-study} are refused unless \code{--engine-mdo-optimize} is set (\code{:3813-3829}).
  \item \textbf{Branch, first match wins:} \code{--requirements} goes to \code{_run_requirements} (\code{:3830}). \code{--engine-mdo-optimize} goes to \code{_run_engine_mdo_optimize} (\code{:3832}). \code{--engine-mdo} goes to \code{_run_engine_mdo} (\code{:3834}). Anything else runs the traditional pipeline of Section~\ref{sec:df-pipeline}.
\end{enumerate}

Figure~\ref{fig:df-pipeline} shows the whole tree: the three entry dispatchers, the job path on the left, the traditional pipeline in the centre in execution order, and the MDO branch on the right. The cross-cutting groups X01 to X08 sit beside the pipeline. They are not steps in it.

\begin{figure}[p]
\centering
\resizebox{\linewidth}{!}{%
\begin{tikzpicture}[
  font=\sffamily\scriptsize,
  >={Stealth[length=1.8mm]},
  bx/.style={draw=#1, rounded corners=2pt, align=center, inner sep=2.5pt, minimum height=7mm},
  side/.style={bx=teal, fill=white, text width=44mm},
  entry/.style={bx=ink, fill=ink!8, text width=44mm},
  stage/.style={bx=teal, fill=pale, text width=50mm},
  mdo/.style={bx=teal, fill=teal!12, text width=44mm},
  xc/.style={bx=muted, dashed, fill=white, text width=44mm, minimum height=6mm},
  note/.style={bx=warn, fill=warnpale, text width=44mm, align=left},
  decision/.style={diamond, aspect=3.2, draw=warn, fill=warnpale, align=center, inner sep=0pt, font=\sffamily\tiny},
  flow/.style={->, draw=ink, thick},
  opt/.style={->, draw=teal, thick, dashed},
  lbl/.style={font=\sffamily\tiny, text=muted, align=center},
  hdr/.style={font=\sffamily\scriptsize\bfseries, text=ink},
  ex/.style={text=warn},
]
% ---------------- column headers
\node[hdr] at (0,1.0) {Job path (\dftt{lrekit} verbs)};
\node[hdr] at (6.0,1.0) {Flag runner (\dftt{run\_nozzle.main})};
\node[hdr] at (12.0,1.0) {Export verb, legacy, MDO};
% ---------------- entry row
\node[entry] (E1) at (0,0) {\textbf{lrekit} $\to$ \dftt{lrekit.cli.main}};
\node[entry, text width=50mm] (E2) at (6.0,0) {\textbf{raosim}, \textbf{RaoRocketSim}, \dftt{python -m}\\$\to$ \dftt{raosim.cli.main}};
\node[entry] (E3) at (12.0,0) {\textbf{RaoRocketSimLegacy} $\to$ \dftt{main.py} (X03): sweeps, benchmarks; \dftt{builds/vNNN\_*/}};
% ---------------- left: job path
\node[side] (J0) at (0,-1.5) {\dftt{lrekit.job\_cli.main};\\\dftt{catalog} prints and exits};
\node[side] (J1) at (0,-3.0) {\textbf{S02} load and resolve the job (\dftt{mdo.job}; \dftt{expand\_architectures} in S14 \dftt{mdo.execute})};
\node[side] (J2) at (0,-4.5) {\textbf{S13} ANALYZE or OPTIMIZE (\dftt{solve\_min\_mass}); engine state};
\node[side] (J3) at (0,-5.85) {\textbf{S14} \dftt{reevaluate} $\to$ \dftt{design\_nozzle\_v2}; Rao contour check; verdict};
\node[side] (J4) at (0,-7.5) {\textbf{S14} run record \dftt{output\_dir/}$\langle$uuid$\rangle$\dftt{/}: \dftt{manifest.json} (SHA-256), \dftt{job.json}, \dftt{summary.json}, \dftt{parts/}};
% ---------------- centre: runner
\node[stage] (R0) at (6.0,-1.5) {\textbf{S01} \dftt{run\_nozzle.main}: 330 flags, \dftt{@}argfiles, menu, shared input validation};
\node[decision] (D) at (6.0,-3.0) {MDO or requirements\\flag, or menu 2/3?};
\node[stage] (T1) at (6.0,-4.5) {\textbf{S01/S02 set-up}: defaults, flag cross-checks, release evidence, interview, wall material};
\node[stage] (T2) at (6.0,-5.9) {\textbf{S03} thermochemistry: \dftt{cea.resolve\_thermochemistry}};
\node[stage] (T3) at (6.0,-7.3) {\textbf{S04} throat sizing and performance; \dftt{-{}-out} created};
\node[stage] (T4) at (6.0,-8.7) {\textbf{S05} B\'ezier/TOP contour (default), chamber, gates $\to$ \dftt{contour.csv}, \dftt{profile.png} {\color{warn}(gate fails: exit 2)}};
\node[stage] (T5) at (6.0,-10.1) {\textbf{S12} early figures: separation, BDE diagrams, flow field, GIFs};
\node[stage] (T6) at (6.0,-11.5) {\textbf{S07} regen wall sizing, Bartz/Sieder--Tate cooling, wall stress};
\node[stage] (T7) at (6.0,-12.9) {\textbf{S08} pintle injector, optional spray/$c^*$ loop $\to$ \dftt{pintle.json} {\color{warn}(rejected: exit 2)}};
\node[stage] (T8) at (6.0,-14.3) {\textbf{S09} electric pumps $\to$ \dftt{pump*.json}; S11 pump CAD; S12 GIF {\color{warn}(gates fail: exit 2)}};
\node[stage] (T9) at (6.0,-15.7) {\textbf{S10} interface sync, throttle map; \textbf{S11} \dftt{pintle/}; \textbf{S12} injector figures};
\node[stage] (T10) at (6.0,-17.1) {\textbf{S11} core CAD: wall/jacket STL and STEP, regen, assembly};
\node[stage] (T11) at (6.0,-18.5) {\textbf{S14} \dftt{summary.json} (\dftt{hardware\_qualified} False); exit 0};
% ---------------- right: export, MDO, S06
\node[side] (OF) at (12.0,-1.75) {\textbf{X02} OpenFOAM VOF case + manifest; OpenFOAM never run};
\node[mdo] (M1) at (12.0,-3.0) {\textbf{S13} requirements, MDO evaluate or MDO optimize (\dftt{mdo.engine}, \dftt{mdo.nlp}; JAX x64)};
\node[mdo] (M2) at (12.0,-4.6) {\textbf{S14} \dftt{reevaluate} $\to$ \dftt{design\_nozzle\_v2} $\to$ v2 files in \dftt{-{}-out}};
\node[mdo] (M3) at (12.0,-6.1) {Console tables; exit code = verdict (0 to 4)};
\node[mdo] (S6) at (12.0,-8.7) {\textbf{S06} (opt-in) \dftt{solve\_rao\_bvp}: MOC/BDE net; \dftt{-{}-backend jax} uses X01};
\node[note] (N1) at (12.0,-12.2) {\textbf{Early returns, traditional path}\\[1pt]
\textbullet\ set-up: tags and material list exit 0; parser errors exit 2\\
\textbullet\ S05: contour or chamber gate fails; no files, exit 2\\
\textbullet\ S08: injector rejected; \dftt{summary.json}, exit 2\\
\textbullet\ S09 to S11: pump or injector gates, a requested pump or pintle CAD package, or the throttle map fail; \dftt{summary.json}, exit 2\\
\textbullet\ B\'ezier chart lookup with $\varepsilon\notin[5,50]$ or length $\notin[60,100]\,\%$: uncaught \dftt{ValueError}};
\node[note, draw=ink, fill=white] (N2) at (12.0,-16.65) {\textbf{Legend}\\[1pt]
solid arrow: default flow\\
dashed arrow: flag-gated branch\\
shaded box: stage, in execution order\\
dashed grey box: cross-cutting group\\
S12 figures and S11 pump CAD run inside the block they follow};
% ---------------- left: cross-cutting panel
\node[hdr] at (0,-8.7) {Cross-cutting groups (not steps)};
\node[xc] (X1) at (0,-9.3) {\textbf{X01} JAX core: LM solve for S06; primitives for S13};
\node[xc] (X2) at (0,-10.62) {\textbf{X02} spray parcels, OpenFOAM export; S08 uses handoff types only};
\node[xc] (X3) at (0,-11.94) {\textbf{X03} legacy toolbox CLI (\dftt{main.py})};
\node[xc] (X4) at (0,-13.26) {\textbf{X04} vendored NASA/JHU MOC reference (test oracle)};
\node[xc] (X5) at (0,-14.58) {\textbf{X05} 15 standalone research scripts};
\node[xc] (X6) at (0,-15.9) {\textbf{X06} knowledge graph, sources, packaged data};
\node[xc] (X7) at (0,-17.22) {\textbf{X07} tests and \dftt{examples/cli/*.args}};
\node[xc] (X8) at (0,-18.54) {\textbf{X08} 8 modules reachable only from tests or nothing};
% ---------------- arrows
\draw[flow] (E1) -- node[lbl, above] {other} node[lbl, below] {argv} (E2);
\draw[flow] (E1) -- node[lbl, right, align=left] {guided $|$ validate $|$\\run $|$ catalog} (J0);
\draw[flow] (J0) -- (J1);
\draw[flow] (J1) -- (J2);
\draw[flow] (J2) -- (J3);
\draw[flow] (J3) -- (J4);
\draw[flow] (E2) -- (R0);
\draw[opt] (E2.east) -- ++(0.45,0) |- (OF.west);
\node[lbl, anchor=west] at (9.15,-0.92) {\dftt{export-openfoam-spray}};
\draw[flow] (R0) -- (D);
\draw[flow] (D) -- node[lbl, above] {yes} (M1);
\draw[flow] (D) -- node[lbl, right] {no: traditional} (T1);
\draw[opt] (M1) -- node[lbl, left] {\dftt{-{}-mdo-export}} (M2);
\draw[flow] (M2) -- (M3);
\draw[flow] (M1.east) -- ++(0.3,0) |- (M3.east);
\draw[flow] (T1) -- (T2);
\draw[flow] (T2) -- (T3);
\draw[flow] (T3) -- (T4);
\draw[opt, <->] (T4) -- node[lbl, above] {\dftt{rao-bvp}} (S6);
\draw[flow] (T4) -- (T5);
\draw[flow] (T5) -- (T6);
\draw[flow] (T6) -- (T7);
\draw[flow] (T7) -- (T8);
\draw[flow] (T8) -- (T9);
\draw[flow] (T9) -- (T10);
\draw[flow] (T10) -- (T11);
\end{tikzpicture}}
\caption{Execution order of the LREKit command line (working tree, 5 November 2026). Centre: the traditional pipeline of run\_nozzle.main, S01 to S14 top to bottom. Right: the export verb, the legacy toolbox and the early-exit MDO branch. Left: the job path of lrekit run with its run record, and the cross-cutting groups.}\label{fig:df-pipeline}
\end{figure}

% ==============================================================================
\section{The traditional pipeline, stage by stage}\label{sec:df-pipeline}

This is the path for \code{lrekit}~\code{[flags]}, for bare runs that choose menu item~1, and for the shipped \code{@examples/cli/*.args} files. The stages run strictly in the order below. Figures (\leavevmode\stagetag{S12}) and some CAD (\leavevmode\stagetag{S11}) are interleaved with the stage blocks they depend on.

\subsection{S01 and S02: defaults, cross-checks, release evidence and interview}\label{sec:df-s01s02}
\begin{itemize}
  \item \code{--complete-package}, or a bare run, calls \code{_apply_complete_package_defaults} (\code{:607}, called at \code{:3836}). It sets LOX/RP-1, \code{--injector}~\code{pintle}, \code{--electric-pump} and the starter tank pressures, and records what it filled in \code{args._complete_package_default_notes}.
  \item \code{_apply_wall_sizing_mode} (\code{:920}, called at \code{:3842}) maps \code{--wall-sizing} and \code{--size-wall} onto \code{args.regen} and \code{args.size_wall}.
  \item The mutual-exclusion and domain checks run at \code{:3843-4005}. For example, frozen variable-$c_p$ expansion is B\'ezier-only and excludes regen and thermal runs, \code{--regen-brep} requires \code{--regen} and \code{--cad}, \code{--flowfield} and \code{--animate} require \code{rao-bvp}, and the spray flags require \code{--spray-cstar-coupling}. Bounds checks cover $\gamma$, $L^*$, contraction ratio, shoulder and throat radii, flange and bolts, wall thickness and liner thickness.
  \item The cooling options (coolant property backend, hydraulic network, ports, radiation, boiling and critical heat flux) are collected into \code{args._cooling_options} (\code{:4013-4028}).
  \item \code{--tags} (\code{print_tags}, \code{:139}) and \code{--list-materials} (\code{materials.material_table}, \code{:4038-4042}) print and exit 0.
  \item \textbf{Release evidence} (\code{:4044-4062}). With \code{--release-evidence-manifest}, the runner calls \code{release_readiness.load_evidence_manifest}. Otherwise it calls \code{evaluate_release_readiness(target="engine")}. The report is stored as \code{args._release_readiness}, and \code{--require-release-evidence} turns an incomplete report into a parser error.
  \item The interactive interview \code{_interactive} (\code{:957}, called at \code{:4073}) runs when the run is bare or uses \code{-i}. The wall-sizing mode and the common validation are then applied again.
  \item \code{_ensure_pyplot} (\code{:324}, called at \code{:4077}) imports Matplotlib lazily. It uses the Agg backend unless \code{--show}, a terminal or an interactive run asks for windows.
  \item The coolant inlet temperature defaults by fluid (\code{:4079-4090}). \code{materials.get_material(--material)} sets the wall conductivity and temperature limit (\code{:4096-4113}). \code{--size-wall} without \code{--material} returns 2.
\end{itemize}

\subsection{S03: propellant and thermochemistry}\label{sec:df-s03}
\begin{itemize}
  \item The propellant comes from \code{--propellant} or from \code{--oxidizer} and \code{--fuel}. The oxidizer and fuel names for the injector feeds are split at \code{:4138-4147}. With \code{--regen}, \code{--spray-cstar-coupling} is refused unless the coolant is the cycle fuel (\code{:4148-4157}).
  \item \code{--mixture-ratio} defaults to the propellant table's nominal O/F (\code{:4165-4182}).
  \item \code{cea.resolve_thermochemistry} (\code{cea.py:211}, called at \code{:4202}) returns the \code{Propellant} record ($\gamma$, molecular weight, chamber temperature, $c^*$, $\eta_{c^*}$, $\eta_{C_F}$). With \code{--thermo-mode}~\code{constant-gamma} it uses the built-in table in \code{propellants.py}. With \code{--thermo-mode}~\code{cea} it uses a RocketCEA frozen chamber snapshot. Without a named propellant, \code{propellants.custom_propellant} builds a placeholder gas (\code{:4184-4199}). Efficiency or $\gamma$ overrides rebuild a custom propellant (\code{:4213-4245}).
  \item \code{--nozzle-expansion-model}~\code{frozen-variable-cp} loads a provenance-bound $c_p(T)$ table through \code{frozen_flow.load_frozen_gas_table} (called at \code{:4257}) and takes $\gamma$ from it.
  \item \code{--cd-target} sets \code{args.ru_factor} through \code{throat_geometry.upstream_radius_ratio_for_discharge_coefficient} (called at \code{:4292}).
\end{itemize}

\subsection{S04: throat sizing and performance}\label{sec:df-s04}
\begin{itemize}
  \item The ambient pressure is $p_a = (p_a/p_0)\,p_c$. \code{--target-thrust} sets \code{args.rt} through \code{design.throat_radius_for_target_thrust} (\code{design.py:1222}, called at \code{:4311}).
  \item \code{engine.compute_engine_performance} (\code{engine.py:69}, called at \code{:4319}) returns $C_F$, $c^*$, $I_{sp}$, $\dot m$, the exit Mach number and the exit pressure. The total mass flow is split into fuel and oxidizer flows by O/F (\code{:4330-4333}).
  \item \code{_apply_split_injector_pressure_model} (\code{:680}, called at \code{:4339}) sets the fuel and oxidizer injector $\Delta p$ as fractions of $p_c$. When the coolant is the cycle fuel, it also sets the regen jacket outlet boundary. \code{_validate_pump_args} (\code{:796}, called at \code{:4340}) follows.
  \item The output directory is created only now (\code{args.out.mkdir}, \code{:4342}). The build plan and the performance panel print at \code{:4345-4440}.
\end{itemize}

\subsection{S05 and S06: nozzle contour and chamber}\label{sec:df-s05s06}
\begin{itemize}
  \item \code{_solve(args)} (\code{:1464}, called at \code{:4452}) selects the contour method. With the default \code{bezier}, it calls \code{nozzle_geometry.bell_nozzle_contour} with angles from the SP-125 Figure~4-14 chart held in \code{rao_top_chart.py}, adds reliability metadata when frozen flow is active, and wraps the result in \code{_TrustedTopSolution} (\code{:1403}).
  \item With \code{rao-bvp}, it builds a \code{RaoSolverConfig} with \code{wall_method="bde"} and calls \code{rao_variational.solve_rao_bvp} (\code{rao_variational.py:5901}). This path uses the MOC unit processes (\code{moc.py}, \code{moc_bounded.py}), the NASA/JHU BDE port (\code{nasa_moc.py}), the transonic kernel, the residuals, the topology checks and the wall splines. \code{--backend}~\code{jax} uses \code{raosim.jax.api.least_squares_jax} (called at \code{rao_variational.py:6069}) and \code{jax.theta_b_solve.least_squares_jax_theta_b} (called at \code{:6029} of the same file). \code{--backend}~\code{numpy} uses SciPy \code{least_squares}.
  \item The reliability gates run at \code{:4456-4840}: the $2\times10^{-3}$ residual gate, mass and length closure, MOC compatibility, the valid region and thrust sanity. If the export gate fails, the run returns 2 with no artifacts (\code{:4834}). The overrides are \code{--allow-unconverged} on \code{rao-bvp} and \code{--allow-chart-extrapolation} on B\'ezier (\code{:4817-4820}).
  \item Chamber (\code{:4841-4896}): \code{ThroatGeometrySpec}, \code{chamber_geometry.auto_shoulder_factor} (called at \code{:4861}) when \code{--shoulder-sizing}~\code{auto}, \code{chamber_contour} with $L^*$-exact volume closure (\code{:4868}), \code{full_engine_contour} (\code{:4876}) and \code{failed_thrust_chamber_geometry_checks} (\code{:4880}). A failed check returns 2 with no artifacts.
  \item The \code{summary} dictionary is assembled at \code{:4899-5040}, including \code{model_registry.audit_model_registry} and \code{model_provenance_dict} (\code{:4967-4974}).
  \item Artifacts: \code{contour.csv} and \code{profile.png} (\code{:5042-5052}), plus \code{contour_moc_crossings.csv} when crossings were sampled (\code{:5053-5062}).
\end{itemize}

\caution{In the working tree, \code{raosim/nozzle_geometry.py} (modified, uncommitted) and the new, untracked \code{raosim/rao_top_chart.py} add a strict chart-domain gate. \code{lookup_angles} (\code{nozzle_geometry.py:47}) raises \code{ValueError} when $\varepsilon$ lies outside $[5,50]$ or the bell length outside $[60,100]\,\%$. \code{_solve} passes no explicit angles, so the exception propagates out of \code{main} as a traceback. It is raised before the \code{--allow-chart-extrapolation} override at \code{:4817-4820} is reached, so that flag cannot rescue an out-of-domain request. On 5 November 2026 the shipped \code{examples/cli/test_13kn_sealevel_regen_allstl.args} ($\varepsilon=4.524$, default B\'ezier) failed this way at \code{run_nozzle.py:4452}. \code{test_13kn_movable_pintle.args} uses the same $\varepsilon$ on the same path. The \code{rao-bvp} example ran to exit 0.}

\subsection{S12: early figures}\label{sec:df-s12early}
These run in best-effort blocks: exceptions are caught and reported as skips.
\begin{itemize}
  \item \code{plotting.plot_separation_on_contour} (called at \code{:5074}), using \code{separation.check_separation}, writes \code{separation_contour.png}.
  \item When \code{sol.construction_diagnostics["bde_artifacts"]} exists (\code{rao-bvp} only), \code{moc_diagrams} writes \code{kernel_fan.png}, \code{rao_topology.png}, \code{bde_mesh.png}, \code{bde_diagnostics.png} and \code{bde_integrity.png} (\code{:5091-5122}).
  \item \code{--flowfield} calls \code{flow_viz.plot_flowfield} (\code{:5130}) and writes \code{flowfield.png}. \code{--animate} calls \code{animate_moc_march} and/or \code{animate_particles} (\code{:5156}, \code{:5167}) and writes \code{anim_moc_march.gif} and \code{anim_particles.gif}.
\end{itemize}

\subsection{S07: regen wall sizing, cooling and wall structure}\label{sec:df-s07}
\begin{itemize}
  \item \code{--regen}~\code{--size-wall} (\code{:5177-5359}). \code{thermal_design.joint_wall_channel_design} (called at \code{:5191}) sizes the hot-wall thickness, channel count, width and height against the thermal, SP-125 eq.~4-31 stress, fatigue and pressure-drop limits, using the cycle coolant flow. \code{size_wall_profile} (called at \code{:5302}) then sizes station-wise hot-wall and jacket thicknesses into a \code{regen_profile.RegenWallProfile} stored as \code{args._wall_profile}. The run sets \code{args.thermal=True} and fills \code{summary["wall_sizing"]} and \code{summary["wall_profile"]}.
  \item \code{--regen}~\code{--auto-size} calls \code{thermal_design.size_cooling_channels} (called at \code{:5367}) and writes \code{summary["channel_sizing"]}.
  \item \code{--regen}~\code{--thermal} (\code{:5409-5616}) runs \code{physics.bartz_heat_flux} and \code{physics.regenerative_cooling_analysis} (called at \code{:5439}). The latter is a Sieder--Tate coolant march, with the property, hydraulic-network, radiation and boiling options supplied by \code{thermofluids.py}. Results go to \code{summary["cooling"]}. Figures: \code{cooling_profile.png}, \code{channel_march.png} and \code{cross_section.png} (from \code{physics.wall_cross_section_field}). With catalog elastic properties, \code{physics.coaxial_shell_wall_stress_profile} (called at \code{:5572}) fills \code{summary["structural"]} and \code{structural_life.png}.
\end{itemize}

\subsection{S08: pintle injector and the optional spray/c* fixed point}\label{sec:df-s08}
\begin{itemize}
  \item The injector is evaluated after cooling, so that its fuel feed state is the regen outlet state.
  \item The \code{InjectorSpec} is assembled at \code{:5701-5849} from the propellant feed, pintle geometry, mechanical, movable-pintle and feed-system specifications. \code{_feed_system_spec_from_args} is at \code{:852}.
  \item Default: the closure \code{evaluate_cli_injector} (\code{:5903}) calls \code{injector.evaluate_pintle_injector} (\code{injector.py:3457}, call at \code{:5912}). It is invoked at \code{:6182}. \code{movable_pintle.py} handles \code{son_continuous_movable}.
  \item With \code{--spray-cstar-coupling}, \code{spray_coupling.solve_spray_cstar_fixed_point} (\code{spray_coupling.py:210}, called at \code{:5989}) wraps the injector and cooling re-evaluation in a fixed point. The run then recomputes \code{compute_engine_performance} (\code{:5998}) and rewrites the cooling, structural and figure artifacts from the final iterate. The runner builds \code{SprayCStarCouplingSpec} without a \code{source} argument (\code{:5934-5941}), so the default \code{source="legacy_screen"} (\code{spray_coupling.py:32}) applies. The Lagrangian parcel solver in \code{raosim/spray/} is therefore never executed from the CLI.
  \item A rejected injector writes \code{summary.json} and returns 2 (\code{:6188-6193}). Otherwise \code{pintle.json} is written (\code{:6204}).
\end{itemize}

\subsection{S09: electric pump feed}\label{sec:df-s09}
\begin{itemize}
  \item This stage runs only with \code{--electric-pump}, which the complete-package defaults switch on. \code{pumps.size_electric_pumps(inj.feed_system, _pump_spec_from_args(args))} is called at \code{:6212} (\code{pumps.py:3747}; \code{_pump_spec_from_args} at \code{:881}). It writes \code{pump.json}, \code{pump_bom.json} and \code{pump_reference_geometry.json} (\code{:6219-6238}).
  \item The \leavevmode\stagetag{S11} pump CAD runs inside this block. \code{pump_cad.export_pump_package} (called at \code{:6250}) writes \code{pump/}, and \code{--pump-cad-format}~\code{step|both} adds \code{pump_cad_brep.export_pump_brep_package} (\code{:6267}). \code{--pump-visualize} adds \code{flow_viz.animate_pump_particles} (\code{:6340}).
  \item Failed pump gates write \code{summary.json} and return 2 unless \code{--allow-infeasible-pump} is given (\code{:6352-6357}). Failed injector gates do the same unless \code{--allow-infeasible-injector} is given (\code{:6361-6365}). A requested pump CAD package that fails a geometry or export gate also writes \code{summary.json} and returns 2 (\code{:6313-6326}).
\end{itemize}

\subsection{S10 and S11: interface sync, throttle map and injector package}\label{sec:df-s10}
\begin{itemize}
  \item \code{injector_cad.resolve_machined_pintle_layout} (called at \code{:6373}) and \code{interface.resolve_bolted_interface_geometry} (called at \code{:6375}) auto-size the flange outer diameter, the injector face outer diameter and thickness, the bolt count, the bolt circle and the holes. The results are written back into \code{args} and the injector specification (\code{:6369-6432}).
  \item \code{--throttle-map} calls \code{injector.throttle_map} (\code{injector.py:2565}, called at \code{:6442}). An unreachable throttle map writes \code{summary.json} and returns 2 (\code{:6500-6511}).
  \item \code{injector_export.export_pintle_package} (called at \code{:6525}) writes the \code{pintle/} folder: always \code{pintle_parameters.json}, \code{pintle_dimensions.csv}, \code{pintle_schematic.svg} and \code{pintle_cross_section.png}; with machined or auto CAD, STEP bodies through \code{injector_cad} and \code{injector_coaxial_cad} and an \code{injector_manufacturing_report.json}. If the package fails while \code{--injector-cad} is not \code{none}, the run writes \code{summary.json} and returns 2 (\code{:6543-6557}); with \code{--injector-cad}~\code{none} the failure is reported as a skipped diagnostic package. \code{injector_plots.export_all_injector_figures} (called at \code{:6563}) writes \code{injector_*.png}.
  \item The interface screen (\code{:6574-6655}) runs \code{interface.screen_composite_regen_wall} (\code{:6592}) when a sized profile and a cooling result exist, then \code{interface.screen_injector_chamber_interface} (\code{:6611}).
\end{itemize}

\subsection{S11: core CAD}\label{sec:df-s11}
\begin{itemize}
  \item \code{export.export_stl} (called at \code{:6689}) writes \code{wall.stl}, with the flange from \code{--flange-od} and \code{--flange-length}. \code{export.inspect_stl} (\code{:6695}) audits that the mesh is watertight. A sized wall profile adds \code{jacket.stl} through \code{regen_profile.normal_offset_contour} (\code{:6728-6733}).
  \item \code{--cad}~\code{step|ipt|both} calls \code{export.export_step} (\code{:6747}, \code{:6785}), which writes \code{wall.step} and, with a sized profile, \code{jacket.step}. A true B-rep is required unless \code{--no-require-brep} is given.
  \item \code{--regen-brep} and \code{--regen-manifolds} call \code{regen_cad.export_channel_wall_step} (\code{:6820}), which writes \code{regen.step} through CadQuery. \code{--cad}~\code{ipt|both} adds \code{export.package_ipt_request} (\code{:6863}) and \code{wall_ipt_manifest.json}. Native IPT is not written.
  \item \code{--regen} calls \code{regen_geometry.generate_regen_nozzle} (\code{:6888}), which writes \code{regen.stl} and \code{regen_3d.png}. \code{--engine-assembly} calls \code{engine_cad.export_engine_assembly} (\code{:6938}), which assembles \code{engine_assembly.step} from the wall, jacket, pintle and pump STEP files written earlier. A failed assembly gate is re-raised as \code{RuntimeError} (\code{:6951-6955}) and ends the run with a traceback.
\end{itemize}

\subsection{S14: final write}\label{sec:df-s14final}
\code{summary.json} is always written (\code{:6957}) and the artifact list is printed. \code{args._release_readiness.blocked} selects between the ``Physical release BLOCKED'' message (\code{:6965-6970}) and the ``evidence complete'' message. The run returns 0 (\code{:6977}).

% ==============================================================================
\section{Early-exit MDO branches}\label{sec:df-mdo}

The three MDO-family flags leave \code{main} before the traditional set-up (Section~\ref{sec:df-dispatch}, step~7). They share these helpers: \code{_mission_from_mdo_args} (\code{:1998}) builds an \code{mdo.schema.MissionSpec}; \code{_resolve_mdo_of_intent} (\code{:1913}) applies \code{requirements.MixtureRatioMode}; and \code{_mdo_contour_request_error} (\code{:1866}) allows \code{rao-bvp} only together with \code{--mdo-export} and \code{--cad}~\code{none}. JAX is imported lazily with 64-bit floats enabled (\code{:2096}, \code{:2516}).

\begin{description}[style=nextline,leftmargin=1.2em]
  \item[\texttt{-{}-engine-mdo} $\to$ \texttt{\_run\_engine\_mdo} (:2077)] This branch runs \code{mdo.engine.chamber_surfaces_for} and \code{mdo.coolant_htd.require_htd_coverage} first; missing coverage returns \code{INDETERMINATE} (\code{:2130}). A \code{DesignVector} goes to \code{mdo.engine.solve_engine} (\code{mdo/engine.py:282}, called at \code{:2174}). That is one coupled, differentiable evaluation of the grid, properties, propellants, cooling (with \code{jax.thermal}), injector, pintle, pump, mass, structures, envelope and assembly blocks. It reports the constraint table and the RQ1 ablation delta. With \code{--mdo-export}, \code{mdo.postprocess.reevaluate} (\code{postprocess.py:971}, called at \code{:2301}) runs \code{design.design_nozzle_v2} (\code{design.py:565}) and \code{design._write_v2_artifacts} (\code{design.py:2589}) into \code{--out}.
  \item[\texttt{-{}-engine-mdo-optimize} $\to$ \texttt{\_run\_engine\_mdo\_optimize} (:2498)] With \code{--isp-sweep}, \code{mdo.nlp.pareto_frontier} (called at \code{:2630}) traces the mass-$I_{sp}$ frontier, and each converged point is re-evaluated into \code{--out/pareto_NNN/} (\code{:2662-2705}). Otherwise \code{mdo.nlp.solve_min_mass} (called at \code{:2723}) runs SLSQP with exact JAX Jacobians. \code{--mdo-station-study} adds \code{mdo.convergence.station_refinement_study} (\code{:2592}), and \code{--mdo-export} re-evaluates into \code{--out} (\code{:2756}).
  \item[\texttt{-{}-requirements} $\to$ \texttt{\_run\_requirements} (:2360)] This branch builds a \code{requirements.EngineRequirement}. The thrust condition is parsed by \code{_parse_thrust_condition} (\code{:2345}), with \code{atmosphere.pressure} for altitude. \code{requirements.solve_requirement} (called at \code{:2427}) calls \code{mdo.nlp.solve_min_mass}. Output is console only.
\end{description}

{\small
\begin{longtable}{P{3.6cm}P{1.0cm}P{10.8cm}}
\caption{Exit codes of the early-exit branches. The enumeration WorkflowExitCode is defined in lines 28 to 35 of raosim/input\_validation.py; \_run\_requirements returns plain integers with the same meaning (run\_nozzle.py lines 2420 to 2476).}\label{tab:df-exit}\\
\toprule
Verdict & Code & Meaning \\
\midrule
\endfirsthead
\toprule
Verdict & Code & Meaning \\
\midrule
\endhead
\code{MET} & 0 & Candidate meets every enforced margin or requirement. \\
\code{CANDIDATE_VIOLATES} & 1 & A solution exists and at least one constraint or requirement is violated. \\
\code{INVALID_INPUT} & 2 & Input or requirement rejected. \\
\code{INDETERMINATE} & 3 & Model coverage missing (for example, no coolant heat-transfer data). \\
\code{SOLVER_FAILED} & 4 & Solver error or non-convergence. \\
\bottomrule
\end{longtable}}

\code{--mdo-export} writes the v2 package into \code{--out}: \code{thrust_chamber_profile.csv}, \code{thrust_chamber_wall.step} or \code{.stl} with its \code{.cad.json} sidecar, \code{design_report_v2.json}, a \code{pintle/} folder and \code{engine_analysis_snapshot_v1.json} (path from \code{postprocess._snapshot_report_path}, \code{postprocess.py:596}). Without \code{--mdo-export}, these branches write no files.

% ==============================================================================
\section{The job path and its run record}\label{sec:df-job}

\code{lrekit}~\code{run}~\code{JOB.yaml} is the only path that publishes a self-describing, hashed run record. In order:
\begin{enumerate}
  \item \code{_cmd_run} (\code{job_cli.py:209}) loads and builds the specification. \code{mdo.execute.expand_architectures} (\code{execute.py:814}) expands architecture alternatives. More than one candidate goes to \code{execute_engine_job_set} (\code{:856}); one candidate goes to \code{execute_engine_job} (\code{:333}).
  \item \code{RunRecord.start} (\code{run_record.py:33}) creates \code{<output_policy.output_dir>/<uuid4hex>/} and writes \code{manifest.json} (status \code{in_progress}, \code{hardware_release} false) and \code{job.json}.
  \item \code{_execute_resolved_engine_job} (\code{execute.py:463}) solves. ANALYZE uses the job design vector; OPTIMIZE uses \code{mdo.nlp.solve_min_mass} (\code{:559}). The state comes from \code{mdo.state.solve_engine_state} (\code{:598}), and the resolved inputs from \code{mdo.resolved_inputs.resolve_engine_inputs} (\code{:607}). \code{mdo.postprocess.reevaluate} (\code{:640}) then runs \code{design_nozzle_v2}, writing CAD only into \code{parts/} when \code{output_policy.parts} is \code{reference_cad}. An optional post-solve Rao contour validation (\code{mdo.contour_validation}, with \code{rao_existence_scan.RaoConvergenceDomain}, \code{:499-505}) follows, then the statuses and the verdict (\code{_verdict}, \code{:248}).
  \item \code{_attach_part_manifest} (\code{:401}) calls \code{mdo.build.build_engine_parts} and \code{write_engine_build}, which write \code{parts/engine_build_manifest.json}.
  \item \code{RunRecord.finish} (\code{run_record.py:72}) writes each requested report through an atomic temporary-file rename and records its SHA-256 digest in \code{manifest.json}. It also hashes the parts manifest and any CAD file that no part claims.
  \item \code{--out}~\code{FILE} writes the result JSON. The process exit code is the verdict (Table~\ref{tab:df-exit}).
\end{enumerate}

\decision{\textbf{Code wins over the trace here.} The candidate reports are \code{summary}, \code{resolved_inputs}, \code{snapshot}, \code{comparison} and \code{coverage}, but the run record writes only the ones requested. \code{RunRecord.start} requests \code{summary} plus \code{output_policy.reports} (\code{run_record.py:40}), and \code{OutputPolicy} defaults to \code{reports=()}, \code{output_dir=None} and \code{parts="none"} (\code{raosim/mdo/job.py:607-615}). A job document that sets no \code{output_dir} therefore produces no run record at all: \code{RunRecord.start} returns \code{None} (\code{run_record.py:35-36}). With an \code{output_dir} and no \code{reports}, the record holds \code{manifest.json}, \code{job.json} and \code{summary.json}. A report that is requested but not produced is written as an explicit \code{unavailable} stub, and the package is then marked \code{partial}; unless the job sets \code{allow_partial_package}, \code{RunRecord.finish} then raises \code{RuntimeError} (\code{run_record.py:145-146}), which \code{execute_engine_job} converts into a \code{SOLVER_FAILED} result.}

% ==============================================================================
\section{Artifacts written per stage}\label{sec:df-artifacts}

Table~\ref{tab:df-artifacts} lists every file the traditional runner can write into \code{--out}, together with the MDO and job outputs. The paths are relative to the output directory.

{\small
\begin{longtable}{P{1.0cm}P{5.3cm}P{5.1cm}P{3.6cm}}
\caption{Artifacts by stage (traditional runner unless stated). Lines refer to raosim/run\_nozzle.py.}\label{tab:df-artifacts}\\
\toprule
Stage & Function (call site) & Artifact files & Condition \\
\midrule
\endfirsthead
\toprule
Stage & Function (call site) & Artifact files & Condition \\
\midrule
\endhead
\leavevmode\stagetag{S04} & \code{args.out.mkdir} (\code{:4342}) & output directory, default \code{builds/nozzle_run} & always (after throat sizing) \\
\leavevmode\stagetag{S02} & \code{job_cli._cmd_guided}, \code{_cmd_validate} & job YAML or JSON; resolved job JSON & \code{--save}; \code{--resolved-out} \\
\leavevmode\stagetag{S05} & \code{np.savetxt} and Matplotlib (\code{:5042-5052}) & \code{contour.csv}, \code{profile.png} & contour and chamber gates pass \\
\leavevmode\stagetag{S06} & \code{_write_moc_crossing_samples} (\code{:427}, written at \code{:5053-5062}) & \code{contour_moc_crossings.csv} & crossings sampled \\
\leavevmode\stagetag{S12} & \code{plotting.plot_separation_on_contour} (\code{:5074}) & \code{separation_contour.png} & best effort \\
\leavevmode\stagetag{S12} & \code{moc_diagrams.plot_kernel_expansion_fan}, \code{plot_rao_topology}, \code{plot_bde_mesh}, \code{plot_bde_diagnostics}, \code{plot_bde_integrity} (\code{:5091-5122}) & \code{kernel_fan.png}, \code{rao_topology.png}, \code{bde_mesh.png}, \code{bde_diagnostics.png}, \code{bde_integrity.png} & \code{rao-bvp} only \\
\leavevmode\stagetag{S12} & \code{flow_viz.plot_flowfield} (\code{:5130}); \code{animate_moc_march}, \code{animate_particles} (\code{:5156}, \code{:5167}) & \code{flowfield.png}; \code{anim_moc_march.gif}, \code{anim_particles.gif} & \code{--flowfield}; \code{--animate} (both \code{rao-bvp}) \\
\leavevmode\stagetag{S07} \leavevmode\stagetag{S12} & \code{plotting.plot_cooling_profile}, \code{plot_coolant_channel_march}, \code{plot_channel_cross_section}, \code{plot_structural_life_dashboard} (\code{:5499-5616}; rewritten at \code{:6091-6151} after the spray fixed point) & \code{cooling_profile.png}, \code{channel_march.png}, \code{cross_section.png}, \code{structural_life.png} & \code{--regen}~\code{--thermal}; structural figure needs catalog elastic data \\
\leavevmode\stagetag{S08} & \code{inj.to_dict()} (\code{:6204}) & \code{pintle.json} & injector evaluated \\
\leavevmode\stagetag{S09} & \code{pumps.size_electric_pumps} (\code{:6212}) & \code{pump.json}, \code{pump_bom.json}, \code{pump_reference_geometry.json} & \code{--electric-pump} \\
\leavevmode\stagetag{S11} & \code{pump_cad.export_pump_package} (\code{:6250}); \code{pump_cad_brep.export_pump_brep_package} (\code{:6267}) & \code{pump/} (\code{pump_parameters.json}, \code{pump_dimensions.csv}, \code{pump_reference_assembly.stl}, \code{pump_reference_assembly_layout.json}); \code{pump/pump_brep/*.step}, \code{pump_cad_units.json} & pump CAD requested; STEP with \code{--pump-cad-format}~\code{step|both} \\
\leavevmode\stagetag{S12} & \code{flow_viz.animate_pump_particles} (\code{:6340}) & \code{pump_particles.gif} & \code{--pump-visualize} \\
\leavevmode\stagetag{S11} & \code{injector_export.export_pintle_package} (\code{:6525}) & \code{pintle/} (\code{pintle_parameters.json}, \code{pintle_dimensions.csv}, \code{pintle_schematic.svg}, \code{pintle_cross_section.png}; STEP bodies and \code{injector_manufacturing_report.json} with CAD) & injector present \\
\leavevmode\stagetag{S12} & \code{injector_plots.export_all_injector_figures} (\code{:6563}) & \code{injector_*.png} & injector present \\
\leavevmode\stagetag{S11} & \code{export.export_stl} (\code{:6689}, \code{:6732}) & \code{wall.stl}; \code{jacket.stl} & always; sized profile \\
\leavevmode\stagetag{S11} & \code{export.export_step} (\code{:6747}, \code{:6785}) & \code{wall.step}; \code{jacket.step} & \code{--cad}~\code{step|ipt|both} \\
\leavevmode\stagetag{S11} & \code{regen_cad.export_channel_wall_step} (\code{:6820}) & \code{regen.step} & \code{--regen-brep} or \code{--regen-manifolds} \\
\leavevmode\stagetag{S11} & \code{export.package_ipt_request} (\code{:6863}) & \code{wall_ipt_manifest.json} & \code{--cad}~\code{ipt|both} \\
\leavevmode\stagetag{S11} & \code{regen_geometry.generate_regen_nozzle} (\code{:6888}) & \code{regen.stl}, \code{regen_3d.png} & \code{--regen} \\
\leavevmode\stagetag{S11} & \code{engine_cad.export_engine_assembly} (\code{:6938}) & \code{engine_assembly.step} & \code{--engine-assembly} \\
\leavevmode\stagetag{S14} & \code{json.dumps(summary)} (\code{:6957}) & \code{summary.json} & always on exit 0; also written just before the exit-2 returns at \code{:6192}, \code{:6323}, \code{:6356}, \code{:6364}, \code{:6507}, \code{:6554} \\
\leavevmode\stagetag{S14} & \code{mdo.postprocess.reevaluate} $\to$ \code{design._write_v2_artifacts} & \code{thrust_chamber_profile.csv}, \code{thrust_chamber_wall.step/.stl}, \code{.cad.json}, \code{design_report_v2.json}, \code{pintle/}, \code{engine_analysis_snapshot_v1.json} & MDO branches with \code{--mdo-export}; \code{pareto_NNN/} per sweep point \\
\leavevmode\stagetag{S14} & \code{RunRecord.start}, \code{finish}; \code{mdo.build.write_engine_build} & \code{<uuid4hex>/manifest.json}, \code{job.json}, \code{summary.json}, requested reports; \code{parts/engine_build_manifest.json} & \code{lrekit}~\code{run} with \code{output_dir}; parts with \code{reference_cad} \\
X02 & \code{openfoam.case.write_openfoam_case} & OpenFOAM case tree, \code{raosim_openfoam_manifest.json} & \code{export-openfoam-spray} \\
X03 & \code{build_log.create_build_dir}, \code{write_metadata} (\code{main.py:407}, \code{:741}, \code{:929}, \code{:1152}) & \code{builds/vNNN_YYYYMMDD_HHMMSS/metadata.txt} and CSV/STL/JSON outputs & \code{RaoRocketSimLegacy} only \\
\bottomrule
\end{longtable}}

The runner itself never imports \code{raosim.build_log}. Versioned \code{builds/vNNN_*} folders therefore come only from the legacy toolbox, mostly when the test suite drives it. On 5 November 2026 the test run created \code{builds/v455_20261105_070346} to \code{builds/v461_20261105_071738}. Each holds one \code{cli_*.csv} and a \code{metadata.txt} whose mode line reads \code{batch}.

% ==============================================================================
\section{Import closure}\label{sec:df-closure}

An AST import closure covers the 161 in-scope modules: \code{raosim/**} (including \code{jax}, \code{mdo}, \code{spray}, \code{openfoam} and \code{data/tools}), \code{lrekit/*}, \code{main.py} and \code{scripts/*}. It follows imports nested inside functions, relative imports, \code{import_module} string literals, parent-package implication and the \code{lrekit.__path__} alias. It reads files only and executes nothing.

{\small
\begin{longtable}{P{4.0cm}P{1.0cm}P{10.4cm}}
\caption{Static reachability of the 161 modules from the console entry points.}\label{tab:df-closure}\\
\toprule
Category & Count & Modules \\
\midrule
\endfirsthead
\toprule
Category & Count & Modules \\
\midrule
\endhead
Reachable from \code{lrekit.cli}, excluding the legacy \code{main.py} edge & 125 & All S01 to S14 modules; \code{raosim/jax}: 8 of 13 (\code{__init__}, \code{api}, \code{assembly}, \code{moc_kernel}, \code{sensitivities}, \code{theta_b_solve} through \code{rao_variational}; \code{primitives} and \code{thermal} through \code{mdo}); \code{raosim/mdo}: all 34; \code{raosim/spray}: all 14, through \code{design.py} $\to$ \code{spray_coupling.py} $\to$ \code{spray.handoff}, used for types and handoff contracts only; \code{raosim/openfoam}: all 3, through the export verb; \code{raosim/legacy_io.py}, through \code{plotting.py}. \\
Reachable only through \code{RaoRocketSimLegacy} & 7 & \code{main.py}, \code{altitude_performance}, \code{benchmarks}, \code{build_log}, \code{conical}, \code{nozzle_comparison}, \code{trade_study} (group X03). \\
Not reachable from any CLI (excluding \code{scripts/}) & 11 & Entry wrapper \code{raosim/__main__.py} (run as \code{python}~\code{-m}~\code{raosim}, not imported). Standalone tools \code{raosim/data/tools/acquire_sources.py} and \code{verify_collection.py} (X06). Library modules (X08): \code{flange_cad} (only \code{tests/test_flange_cad.py}), \code{trajectory} (only \code{tests/test_trajectory.py}), \code{jax/bvp} (\code{tests/test_jax_bvp.py}), \code{jax/design_opt} (\code{tests/test_jax_thermal_design_opt.py}), \code{jax/pack} (\code{tests/test_jax_bvp.py}), \code{jax/residuals} (tests only), \code{thermal_correlations} (no Python importer), \code{jax/reference} (no importer). \\
Standalone scripts & 18 & All \code{scripts/*.py}: the compatibility wrapper \code{scripts/run_nozzle.py} (S01), two NASA/JHU comparison scripts (X04) and 15 research and generator scripts (X05). \\
\midrule
Total & 161 & \\
\bottomrule
\end{longtable}}

% ==============================================================================
\section{Where the documentation disagrees with the code}\label{sec:df-drift}

Table~\ref{tab:df-drift} starts from the code-versus-document disagreements recorded in the triage taxonomy. Each row was re-checked against the working tree on 5 November 2026. Rows the 2026-11-05 edits already fixed are kept and marked as fixed, so that the earlier claim can be traced. Paths are the post-reorganization locations.

{\small
\begin{longtable}{P{0.4cm}P{3.4cm}P{4.4cm}P{5.1cm}P{1.6cm}}
\caption{Documentation versus code, verified 2026-11-05. In every row the code is authoritative.}\label{tab:df-drift}\\
\toprule
\# & Document, location & Document says & Code does & Status \\
\midrule
\endfirsthead
\toprule
\# & Document, location & Document says & Code does & Status \\
\midrule
\endhead
1 & \code{README.md:1217} (Outputs) & ``Normal CLI runs create \code{builds/vNNN_YYYYMMDD_HHMMSS/}'', with design-gate JSON and \code{metadata.txt}. & The runner writes to \code{--out} (default \code{builds/nozzle_run}, \code{:3534}), never imports \code{build_log}, and reports in \code{summary.json}. Versioned folders come only from legacy \code{main.py} (\code{:407}, \code{:741}, \code{:929}, \code{:1152}), and \code{design_report_v2.json} only from \code{design._write_v2_artifacts}. & open \\
2 & \code{README.md:959-1149} (Command-line usage) & Documents the flag runner and the legacy alias only. & It does not mention the job verbs (\code{job_cli.py:29}), \code{export-openfoam-spray} (\code{raosim/cli.py:11}), or \code{--requirements}, \code{--engine-mdo}, \code{--engine-mdo-optimize} and \code{--mdo-export}, all dispatched before the traditional path (\code{:3830-3835}). The layout block (\code{README.md:1328}) names the job verbs in one line. & open \\
3 & \code{README.md:972-973} & ``Bare \code{lrekit} runs enter the same complete-package interactive interview by default.'' & A bare run first prints \code{_MODE_MENU} (\code{:1756}, \code{:3797-3811}). Only choice~1 enters the complete-package interview. & open \\
4 & \code{README.md:1314-1356} (Repository layout) & Before the 2026-11-05 edit: \code{main.py} was the ``Primary interactive/batch CLI''; \code{lrekit/}, \code{mdo/} and others were missing; \code{trajectory.py} was listed as active. & The working-tree README now lists \code{main.py} as the legacy CLI behind \code{RaoRocketSimLegacy}, adds \code{lrekit/} and links this data-flow map. \code{README.md:826} states that the trajectory module is not wired into the CLI. & fixed \\
5 & \code{lrekit/cli.py:3}; \code{lrekit/job_cli.py:7-13} & Job verbs are ``run/validate/catalog''; the catalog lists six topics. & \code{JOB_VERBS} also contains \code{guided} (\code{job_cli.py:29}), and \code{_CATALOG_TOPICS} has eight, adding \code{unavailable} and \code{capabilities} (\code{:31-34}). & open \\
6 & \code{docs/architecture/MDO_REMEDIATION_AND_OUTPUT_CONTRACT.md:151}, \code{:157} & \code{job_cli.py} ``adds three verbs''; the catalog list ends at \code{unavailable}. & Four verbs (\code{guided} as well). The catalog has eight topics, including \code{capabilities} (\code{job_cli.py:31-34}). & open \\
7 & \code{docs/program/LREKIT_HARDWARE_PROGRAM_2026-09-05.md:620-630} & The proposed verb family ``stands'': \code{job}, \code{component}, \code{build}, \code{test}, \code{calibrate}, \code{ui}, and \code{catalog} with topics \code{models}, \code{components}, \code{processes} and \code{coefficients}. & None of these verbs or catalog topics exist; the implemented set is the four job verbs and eight topics above. & plan only \\
8 & \code{main.py:6-11} docstring; \code{main.py:95-99} epilog & \code{lrekit}~\code{--Rt}~\code{25}~\code{--Pc}~\code{60}; \code{lrekit}~\code{--sweep}. & \code{lrekit} now dispatches to the runner (SI \code{--rt} and \code{--pc}, no \code{--sweep}). Those flags work only under \code{RaoRocketSimLegacy}. & open \\
9 & \code{raosim/run_nozzle.py:20-55} docstring; \code{examples/cli/lox_rp1_13kn_copper_316l_full_engine.args:27-30}, \code{test_13kn_sealevel_regen_allstl.args:56-58}, \code{test_13kn_movable_pintle.args:34-35} & \code{--backend}~\code{jax}~\code{--regen} and \code{--max-nfev}~\code{0} ``seed'' examples; argument files set \code{--max-nfev}, \code{--n-control} and \code{--n-kernel} as if a solve runs. The docstring also implies \code{summary.json} needs \code{--regen}. & None of these select \code{--contour-method}~\code{rao-bvp}, so the contour is the default B\'ezier/TOP chart (\code{:2827-2831}) and \code{_solve} ignores the solver flags (\code{:1464-1487}). Only \code{test_13kn_rao_bvp_flowfield.args} runs the BVP. \code{summary.json} is always written (\code{:6957}). & open \\
10 & \code{--allow-chart-extrapolation} help (\code{run_nozzle.py:2832-2836}) & Writes diagnostic B\'ezier artifacts when $\varepsilon$ or the length lies outside the chart domain. & The working-tree \code{lookup_angles} raises \code{ValueError} outside the domain (\code{nozzle_geometry.py:47}, raise at \code{:67}) before the override at \code{:4817-4820}. The flag can no longer produce out-of-domain artifacts. \code{test_13kn_sealevel_regen_allstl.args} ($\varepsilon=4.524$) was observed failing; \code{test_13kn_movable_pintle.args} uses the same $\varepsilon$ on the same path and is expected to fail (not executed). & open (working tree) \\
11 & \code{docs/audits/FEATURE_DATAFLOW_AUDIT_2026-07-11.md:20-70} & The packaged CLI runs \code{design._build_v2_contour}, \code{physics.boundary_layer_displacement}, \code{validation.evaluate_design_gates} and \code{build_log}; legacy adds a trajectory integrator; spray/$c^*$ is blocked with CLI regen; the runner has 4,558 lines. & \code{run_nozzle.main} calls none of the four; they live in \code{design_nozzle_v2}. \code{main.py} does not import \code{trajectory}. Spray with regen is allowed when the coolant is the cycle fuel (\code{:4148-4157}). The runner has 6,981 lines. & historical record \\
12 & \code{raosim/jax/__init__.py:7}; \code{raosim/jax/bvp.py:2-4} & JAX ``is the default for \code{solve_rao_bvp}''; \code{jax/bvp.py} replaces SciPy \code{least_squares}. & The first holds for \code{--backend}, but the default CLI contour never calls \code{solve_rao_bvp}. The live JAX solve is \code{jax/api.least_squares_jax} (\code{rao_variational.py:6069}); \code{jax/bvp.py} is imported only by \code{tests/test_jax_bvp.py}. & open \\
13 & \code{raosim/flange_cad.py:3} (``The gap this closes''); \code{raosim/thermal_correlations.py:1-5} & \code{flange_cad} presents itself as filling the missing chamber-side flange of the bolted joint; \code{thermal_correlations} presents itself as source-scoped correlation choices and applicability reports (Bartz bounds). & Neither is imported by any CLI path. The live flange comes from \code{export.export_step} and \code{interface.resolve_bolted_interface_geometry}, and \code{physics.bartz_heat_flux} does not consult \code{thermal_correlations}. & open \\
14 & \code{docs/subsystems/S08-injector-and-spray-coupling/OPENFOAM_SPRAY_EXPORT.md:88} & \code{lrekit}~\code{export-openfoam-spray}~\code{...}~\code{--case-row}~\code{case_1}. & Matches \code{raosim/cli.py:11} and \code{case.py:1611-1627}. & consistent \\
\bottomrule
\end{longtable}}

% ==============================================================================
\section{Finding things}\label{sec:df-find}

\begin{itemize}
  \item \textbf{Stage table and flow in markdown:} \code{docs/architecture/DATAFLOW.md}, the twin of this chapter.
  \item \textbf{Every code file, its stage and a one-line role:} \code{docs/architecture/CODE_MAP.md}, ordered S01 to S14, then X01 to X08.
  \item \textbf{Verbatim module docstrings and leading comments, by stage:} \code{docs/architecture/CODE_HEADER_COMMENTS.md}.
  \item \textbf{Physics and design basis per stage:} \code{docs/subsystems/Sxx-<slug>/}. Folders exist for S04, S05, S06, S07, S08, S09, S10 and S13. S01, S02, S03, S11, S12 and S14 have no subsystem folder; their contracts are documented in \code{docs/architecture/} (\code{MDO_REMEDIATION_AND_OUTPUT_CONTRACT.md}, \code{PHYSICAL_RELEASE_GATES.md}, \code{MODEL_REGISTRY.md}).
  \item \textbf{Per-stage chapters of this guide:} S01 and S02 in Chapter~\ref{ch:s01-s02}; S03 and S04 in Chapter~\ref{ch:s03-s04}; S05 and S06 in Chapter~\ref{ch:s05-s06}; S07 in Chapter~\ref{ch:s07}; S08 and S09 in Chapter~\ref{ch:s08-s09}; S10 to S12 in Chapter~\ref{ch:s10-s12}; S13 and S14 in Chapter~\ref{ch:s13-s14}; the cross-cutting groups in Chapters~\ref{ch:crosscutting-a} and \ref{ch:crosscutting-b}.
  \item \textbf{To find where a flag acts:} search \code{raosim/run_nozzle.py} for \code{args.<flag_with_underscores>}. Table~\ref{tab:df-flag-families} gives the consuming stage, and Section~\ref{sec:df-pipeline} gives the line range of that stage's block.
\end{itemize}
